% titel: Blatt 7, A5: Differenzierbarkeit
% skript: Kap. 4.1 (Definition der Ableitung)

% Teil I: Werkzeuge, kein Lösungsweg.

Eine Funktion liegt in $C^1(\mathbb R)$, wenn sie überall differenzierbar ist
und ihre Ableitung stetig ist. Die Schwierigkeit liegt meist an
Nahtstellen, wo eine Fallunterscheidung oder eine hebbare Definitionslücke
(z. B. $\sin x / x$ bei $x=0$) auftritt.

\begin{werkzeugkasten}[Definition der Ableitung]
Die Ableitung an der Stelle $x_0$ ist der Grenzwert des Differenzenquotienten,
\begin{equation}
  f'(x_0)=\lim_{x\to x_0}\frac{f(x)-f(x_0)}{x-x_0}
  =\lim_{h\to0}\frac{f(x_0+h)-f(x_0)}{h},
\end{equation}
sofern er existiert. Für $C^1$ muss zusätzlich der \emph{Grenzwert der
Ableitungsfunktion} stetig sein: $\lim_{x\to x_0} f'(x)=f'(x_0)$.

\textit{Skript:} Kap. 4.1, Definition 4.1.1.
\end{werkzeugkasten}

\begin{werkzeugkasten}[Nahtstellen und $\sin x/x$ bei 0]
An Nahtstellen (Fallunterscheidung) separat die Ableitung über die Definition
berechnen und beide Seiten vergleichen. Für Quotienten, die bei $0$ zunächst
undefiniert sind, hilft die Regel von de l'Hôpital, z. B.
$\lim_{x\to0}\frac{\sin x}{x}=1$; die Definition setzt dann den Wert $1$
für $f(0)$ voraus, damit $f$ stetig und differenzierbar wird.

\textit{Skript:} Kap. 4.1 und Kap. auf Grenzwerte.
\end{werkzeugkasten}

\begin{achtung}[Darauf achten]
\begin{itemize}
  \item An der Nahtstelle nicht nur beide Ableitungsformeln vergleichen, sondern
        die Ableitung über die Definition am Punkt selbst prüfen.
  \item Stetigkeit von $f$ und Existenz von $f'$ sind nötig, aber nicht
        hinreichend für $C^1$; zusätzlich die Stetigkeit von $f'$ verlangen.
  \item $\sin x/x$ ist bei $x=0$ zunächst nicht definiert; der Wert dort wird
        durch den Grenzwert festgelegt.
\end{itemize}
\end{achtung}
